\ffigbox[\FBwidth]{%
\caption{\centering Graphe \(G\) avec un 3-cycle \(C\) et un 2-sommet \(x\) puis un 2-sommet \(y\)}\label{fig:c3_plus_2_plus_2}
}{
    \fbox{
        \begin{tikzpicture}[scale=1, main node/.style={circle, draw, fill=blue!20, inner sep=1pt, font=\scriptsize, minimum size=5mm, text=black}]
            
            % construction de C (réorganisé)
            \node[main node, fill=green!20] (c0) at (0, 0) {\(c_0\)};
            \node[main node, fill=green!20] (c1) at (0, 2) {\(c_1\)};
            \node[main node, fill=green!20] (c2) at (2, 0) {\(c_2\)};
            
            \draw[draw=green] (c0) -- (c1);
            \draw[draw=green] (c1) -- (c2);
            \draw[draw=green] (c2) -- (c0);

            % rajout du 2-sommet
            \node[main node, fill=green!20] (x) at (2, 2) {\(x\)};
            \node[main node, fill=blue!20] (y) at (3, 1) {\(y\)};
            
            \draw[draw=green] (x) -- (c0);
            \draw[draw=green] (x) -- (c1);
            \draw[draw=blue] (y) -- (x);
            \draw[draw=blue] (y) -- (c2);

            % ellipse autour de G
            \node[ellipse, draw, thick, fill=blue!20,
                minimum width=2.6cm, minimum height=3.6cm] (Gnode) at (5,1) {\(G \setminus A\)};
            % lignes vers G (au lieu de vers le haut)
            
            \draw[draw=blue, dashed] (y) -- (Gnode);
        \end{tikzpicture}
    }
}